\documentclass{beamer}
\usepackage[latin1]{inputenc}
\usepackage[spanish]{babel}
\usepackage{graphicx}

\title[swarm intelligente]{Creación de diapositivas con Beamer}
\subtitle{Saber mas acerca de este programa}
\author{Lorena Diaz Morales \\ Técnicas experimentales}
\date[15-03-2013]{Marzo - 2013}

\begin{document}

\begin{frame}
 \includegraphics[height=0.2\textheight]{logoullvert.eps}
 \includegraphics[height=0.22\textheight]{facultadmatematicas.eps}
 \titlepage
 
 \scriptsize
 \begin{center}
  Facultad de matemáticas \\
  Universidad de La Laguna \\
 \end{center}

 \end{frame}

\begin{frame}
 \frametitle{Índice}
 \tableofcontents
 \end{frame}
 
\section{Introducción} 
Beamer es una clase \latex para generar transparencias para presentaciones.
Un documento Beamer consiste de una sucesion de marcos(frames).
El programa Beamer significa formalmente un proyector digital y viene del alemán.
Fue desarrollado por Till Tantau.
\subsection{Compilar}
Un documento Beamer se puede compilar directamente con pdflatex y verlo
con acrobat reader. Pero...
Si hay gráficos ,pdflatex funciona bien si los graficos estan en pdflatex
si hay graficos eps, podemos cambiar el formato de todas las imagenes a pdf con 
Acrobat Distiller o Adobe Illustrator, por ejemplo.
Beamer carga ``graphicx'' para el soporte de gráficos.



\end{document}